% Fresh phase-three supplement.  All variable counts include the input.
\section{What the newer constructions do over number fields}
\label{sec:new-constructions-number-fields-phase3}

Three different questions must be separated.  A fixed polynomial
over $\mathbb Q$ has an existential projection over every extension
field.  An equation-combining procedure may preserve the projection
of a specified source system over that extension.  Neither fact
identifies the resulting set with the ring of integers of the
extension.  The following results locate explicit failures of the
first and third assertions, and give an effective repair of the
second in Section~\ref{sec:division-norm-compiler-phase3}.
The counterexamples and algebraic encoding results here do not
depend on accepting the source manuscript's arithmetic theorem.

\subsection{The four literal quartics accept every input over $\mathbb Q(i)$}

Let $P_N$ denote the recorded sum-of-squares quartic with $N$
total coordinates, indexed from zero, where $X_0$ is the input.
Set all witnesses to zero except the following two coordinates:
\[
\begin{array}{r|r|r|r}
N&\text{quadratic residuals}&u\text{ coordinate}&v\text{ coordinate}\\\hline
1315&1272&72&73\\
1339&1296&96&97\\
1367&1322&62&78\\
1391&1346&86&102
\end{array}
\]
The exact restriction in each case is
\begin{equation}\label{eq:quartic-two-squares-restriction-phase3}
                     P_N(T,0,\ldots,u,\ldots,v,\ldots,0)
                        =C_N+T^4+u^2+v^2.
\end{equation}
Here $C_N$ is the constant obtained by the same restriction at
$T=u=v=0$.  For $N=1315,1367$ it includes the term
$1024\kappa^2$, with the fixed integer $\kappa=3^{2^{28}}$.
The full integer constants, coordinate names, and hashes of the
four complete coefficient tables are recorded in
\texttt{quartic-base-change-verification\_20261106\_v1.json}.
Defining $C_N$ by this specialization avoids any ambiguity from
rounding or decimal transcription.

\begin{proposition}\label{prop:quartic-collapse-Qi-phase3}
Each of these four literal quartics has existential input
projection equal to $K$ over every characteristic-zero field
$K$ containing $i$.  In particular none defines
$\mathcal O_{\mathbb Q(i)}$ in $\mathbb Q(i)$.
\end{proposition}
\begin{proof}
For any $T\in K$ put $R=C_N+T^4$ and take
\[
                         u=\frac{1-R}{2},\qquad
                         v=i\frac{1+R}{2}.
\]
Then $u^2+v^2=-R$, so
\eqref{eq:quartic-two-squares-restriction-phase3} vanishes.
Taking $T=1/2$ gives a nonintegral input with witnesses already
in $\mathbb Q(i)$.  This identity is exact for the fixed integer
$\kappa$ as well as before its substitution.
\end{proof}

The independent script
\texttt{check\_quartic\_base\_change\_20261106\_v1.py}
obtains the restriction twice: from the expanded sparse polynomial
and by squaring the independently restricted primitive residuals.
It compares the two results and substitutes the displayed generic
witnesses using exact arithmetic in
$\mathbb Q(i)[T,\kappa]$.  The conclusion is therefore stronger
than the general observation that sums of squares can cancel in a
nonreal field.

\subsection{An exact integral input rejected by the unshifted sources}

The unshifted source uses $\beta=X^2+2$ and guards the scalar
contact expression
\begin{equation}\label{eq:scalar-contact-polynomial-phase3}
       r_\beta=\beta(\beta-1)(\beta-2401)
                         (32\beta+343)(2\beta-1)
\end{equation}
away from zero.  Its later nonzero rational-square rescaling does
not remove these zeros.  The individual factors are present as
nonvanishing guards in the indexed coefficient specifications.
Consequently $X=\sqrt{2399}$, $i$, and $\sqrt{-2}$ are rejected
by the unshifted guarded source, even though each is an algebraic
integer in the corresponding quadratic field.

For the indexed polynomials the failure is particularly exact.
Recall the notation of the finite-algebra specification:
$G$ is its product of guards, $n$ is its algebra rank,
\[
 H(V)=\det(I-VM_{G\lambda})=\sum_{k=0}^n a_kV^k,
 \qquad
 R_E(V)=\operatorname{Tr}\bigl(M_E\operatorname{adj}
                       (I-VM_{G\lambda})\bigr).
\]
Each branch has an index $1\le d\le n$ and uses
$g=a_d\operatorname{Res}(H_d,H_d')$, $L=a_d^2g$,
\[
 \Psi=1+\sum_{j=1}^d a_{d-j}a_d^{2j-1}g^jU^j,
 \qquad
 \widetilde R_E=\sum_{k=0}^{n-1}[V^k]R_E(V)(LU)^{n-1-k}.
\]
Its final residuals are $\Psi$, the high coefficients
$a_{d+1},\ldots,a_n$, and the $\widetilde R_E$.

\begin{proposition}\label{prop:indexed-contact-value-one-phase3}
For each unshifted indexed polynomial with respectively
41, 43, 43, and 45 total coordinates from phases three and two,
the polynomial equals $1$ identically in all remaining
coordinates after imposing any of
\[
                         X^2=2399,\qquad X^2=-1,
                              \qquad X^2=-2.
\]
The assertion holds over every characteristic-zero field.
\end{proposition}
\begin{proof}
The indicated contact factor is zero as a polynomial in the
retained base coordinates, so $G=0$ in every specialized
finite algebra, for every value of every finite parameter.
Thus $M_{G\lambda}=0$, $H=1$, and all $a_k$ with $k\ge1$
vanish.  Hence $g=L=0$ and $\Psi=1$.
Moreover $R_E(V)=\operatorname{Tr}(M_E)$ is constant.
Every branch retains nine quartic point generators, so
$n\ge4^9>1$.  Therefore its reversed trace is zero:
the sole possibly nonzero coefficient has the factor
$(LU)^{n-1}=0$.  Every branch sum of squares is consequently
$1$, and so is their finite product.  This calculation does
not use anisotropy of a sum of squares.
\end{proof}

The contact checker verifies a literal multiplicative path from
the zero factor to the materialized guard, evaluates that factor
in $\mathbb Q[X]/(X^2-a)$, and checks the branch-rank assertion.
Its separate checks of the four quadratic systems follow exact
input-dependent gate equations until an inverse equation becomes
$-1=0$.  The recorded reversible row operations transport this
contradiction to each final quadratic residual system.  Thus a
faithful norm combination of the unshifted residuals still rejects
these inputs.  Over $\mathbb Q(\sqrt{2399})$ the literal sum of
squares is faithful as well, and therefore rejects that integral
input.  Over $\mathbb Q(i)$ the literal quartic instead accepts
every input by Proposition~\ref{prop:quartic-collapse-Qi-phase3};
this is precisely the distinction between a system and its
sum-of-squares encoding.

These finite contact obstructions admit the field-dependent shift
repair of Section~\ref{sec:p3-beta-shift}.  The assertion above
is about the unshifted artifacts and does not apply to the
separately emitted $\beta=X^2+2402$ artifacts.  Removing a finite
list of excluded inputs does not restore the elliptic arithmetic
argument considered next.

\subsection{A totally real quartic obstruction to the finite elliptic anchor}

Put $F=\mathbb Q(\sqrt2)$, $d=3-\sqrt2$ and
$D=d^2=11-6\sqrt2$.  Let
\[
                         E:\ y^2=x^3-Dx,
\]
and let $P\in E(F)$ be the fixed point used in the finite
anchor.  Its particular coordinates are immaterial below;
only $P\in E(F)$ is used.  The anchor excludes equality with
the finitely many abscissas $x([n]P)$, $1\le n\le M$, by
taking the product of their nonzero rational norms.

\begin{proposition}\label{prop:totally-real-anchor-counterexample-phase3}
There is a totally real quartic number field $K$, with
$\sqrt2\notin K$, and a point $Q\in2E(K(\sqrt2))$ of
fixed height such that
\[
              x(Q)\in F,\qquad x(Q)\ne x([n]P)
                       \quad\text{for every }n\in\mathbb Z
\]
whenever the abscissa on the right is defined.  The same fixed
point therefore satisfies every finite anchor exclusion, for
every positive bound $M$.
\end{proposition}
\begin{proof}
Let $s$ be a root of
\[
                         s^4-1696s^2+18432=0
\]
and put $K=\mathbb Q(s)$.  Set $F'=\mathbb Q(\sqrt{2737})$,
where $2737=7\cdot17\cdot23$.  The two values of $s^2$ are
$848\pm16\sqrt{2737}$, both strictly positive.  They have
norm $18432=2\cdot96^2$, which is not a rational square.
Thus neither is a square in $F'$, and $[K:\mathbb Q]=4$.
All four conjugates of $s$ are real, so $K$ is totally real.

If $K$ contained $\sqrt2$, it would equal $F'(\sqrt2)$.
Writing $s=a+b\sqrt2$ with $a,b\in F'$ and using
$s^2\in F'$ gives $ab=0$.  In either case the norm of
$s^2$ from $F'$ to $\mathbb Q$ would be a rational square:
it would be $N(a)^2$ or $4N(b)^2$.  This contradicts the
computed norm.  Hence $\sqrt2\notin K$.

In $K(\sqrt2)$ set
\[
                  b=\frac{s}{2}+\frac{48}{s}\sqrt2,
                          \qquad R=(8,b).
\]
The defining equation for $s$ gives
$b^2=424+48\sqrt2=8^3-8D$, so $R\in E(K(\sqrt2))$.
It has nonzero ordinate.  The doubling formula gives
\begin{align*}
 x(Q)=x([2]R)
   &=\frac{(8^2+D)^2}{4b^2}
     =\frac{5697-900\sqrt2}{1696+192\sqrt2}\\
   &=\frac{312741}{87584}
                          -\frac{40941}{43792}\sqrt2\in F.
\end{align*}
This value is not $0,d$, or $-d$, as is already seen from
its coefficient of $\sqrt2$.  Thus $y(Q)\ne0$.

On the other hand $b\notin F$.  If it belonged to $F$, the
norm of its square would be a square in $\mathbb Q$; but
\[
           N_{F/\mathbb Q}(b^2)=424^2-2\cdot48^2
                           =175168=2^6\cdot7\cdot17\cdot23
\]
is not a square.  The doubling slope is
$(3\cdot8^2-D)/(2b)$.  Since $b^2\in F$, the resulting
ordinate $y(Q)$ is a scalar in $F$ times $b$.  That scalar
is nonzero because $y(Q)\ne0$, so $y(Q)\notin F$.

Every multiple $[n]P$ has coordinates in $F$.  Equality of
its abscissa with $x(Q)$ would imply
$Q=[n]P$ or $Q=-[n]P$, and would put $y(Q)$ in $F$, a
contradiction.  Hence each defined difference
$x(Q)-x([n]P)$ is a nonzero element of $F$ and has nonzero
rational norm.  Every finite product is nonzero.
\end{proof}

The halfpoint is represented in the literal two-coordinate
model by $a_0=8$, $a_1=0$, $b_0=s/2$, $b_1=48/s$ in $K$.
Thus the example applies to the extended point chart itself;
it is not an artifact of a split quadratic algebra.
The script \texttt{check\_real\_anchor\_20261106\_v1.py}
checks all curve and doubling identities by exact arithmetic
in the degree-eight compositum.

The conclusion disproves the extension of the inference
``finite exclusion forces a large integer elliptic label.''
Over the original field that inference depends on a cyclic
group description.  Enlarging the field introduces points
outside that fixed cyclic group.  We have not used this
single point to construct witnesses for every other source
equation at a nonintegral input; the proposition is a
counterexample to an essential arithmetic inference, not a
claim of a complete false positive for the entire source.

\subsection{A domain assertion that does extend to fields with a real place}

One useful component survives more generally.  Let
\[
 G_J(W)=1+\sum_{j=1}^J c_jW^j,\qquad
 c_1=-\frac25,\qquad
 \frac{c_j}{c_{j-1}}=\frac{4j-5}{4j+1}\quad(j\ge2).
\]
All $c_j$ are negative, and induction gives the exact identity
\[
                        G_J(1)=-\frac{4J-1}{2}c_J>0.
\]
Consequently $G_J(w)\ge G_J(1)>0$ for $0\le w\le1$.
At a real embedding of the curve above, $D>0$.  For an
affine double $Q=[2](a,b)$ with the required denominators
and $y(Q)$ nonzero,
\[
                  x(Q)=\left(\frac{a^2+D}{2b}\right)^2>0.
\]
The curve equation then implies $x(Q)>\sqrt D$, and hence
$0<D/x(Q)^2<1$.  Thus
\[
             -\frac{x(Q)}{y(Q)}G_J\!\left(\frac D{x(Q)^2}\right)
                          \ne0.
\]
An embedding is injective, so this proves the nonvanishing
in any number field with a real embedding.  If the formal
quadratic algebra splits, apply the argument to each component;
the source denominator guards ensure both are defined.

The analogous assertion is false over unrestricted number
fields.  Choose an algebraic root $w$ of $G_J$.  It is neither
zero nor one.  Choose $x=d/\sqrt w$ and a nonzero square root
$y$ of $x^3-Dx$.  This gives an algebraic point $Q$ where the
displayed logarithm truncation vanishes.  Adjoin a halfpoint
of $Q$, which exists over a finite extension because the
multiplication-by-two morphism is surjective, and, if desired,
adjoin $i$.  The resulting number field supplies the asserted
failure even on $2E$.  Therefore a final norm combination
repairs conjunctions, but does not automatically justify all
quotient-domain simplifications originally proved over
$\mathbb Q$.

\subsection{No one fixed existential formula works over every number field}

\begin{theorem}\label{thm:no-uniform-existential-integer-rings-phase3}
Let $k_0$ be a number field.  Fix an existential formula over
$k_0$ whose quantifier-free part is a finite union of systems
of polynomial equations and nonvanishing conditions.
Suppose its input projection over $k_0$ contains $\mathbb Z$.
There is a finite extension $K/k_0$ over which its input
projection contains a rational noninteger.  In particular no
one such fixed formula defines $\mathcal O_K$ in every finite
extension $K/k_0$.
\end{theorem}
\begin{proof}
Over an algebraic closure $\overline{k_0}$ the formula's
witness locus is constructible.  Its image in the affine
input line is constructible by Chevalley's theorem
\cite[Tag 00FE]{StacksChevalley}.  A constructible subset of
an affine line over an algebraically closed field is finite
or cofinite.  This image contains $\mathbb Z$, so it is
cofinite.  Choose $q\in\mathbb Q\setminus\mathbb Z$ in it.
The finitely many algebraic witness coordinates belong to
some finite extension $K/k_0$.  The formula therefore holds
at $q$ over $K$.  But a rational algebraic integer is an
integer: if $q=a/b$ in lowest terms satisfies a monic
integer polynomial, clearing denominators shows $b\mid a^m$
and hence $b=1$.  Thus $q\notin\mathcal O_K$.
\end{proof}

This includes any finite number of simultaneous equations,
any finite alternatives, and any finite list of denominator
guards.  Increasing these counts cannot overcome the
uniformity obstruction.  It does not preclude choosing a
different formula for each number field.

There is a useful real refinement when the coefficient field
is $\mathbb Q$, or more generally when $k_0$ has a chosen
real embedding.  The real input projection is semialgebraic
by quantifier elimination and contains all positive and
negative integers.  A semialgebraic subset of the line is
a finite union of points and intervals; it therefore contains
both sufficiently distant tails.  Choose a nonintegral
rational input in one tail.  The transfer principle for real
closed fields supplies real algebraic witnesses.  They,
together with $k_0$, generate a number field $K$ with a real
embedding where that input is accepted.  Quantifier
elimination and this transfer principle are stated, for
example, in \cite[Theorem 2.1]{Basu2014} and
\cite[Corollary 4.5 and Example 4.6]{Plaumann2023}.
The resulting field need not be totally real; only a real
embedding is asserted.

\subsection{What remains open after these checks}

None of these counterexamples proves that some
$\mathcal O_K$ is not Diophantine in $K$.  They prove that
the literal formulas, fixed cyclic arithmetic argument,
or fixed field-independent coefficient choice cannot supply
the desired universal conclusion.  A field-adapted norm form
faithfully combines an already correct finite system; it
does not produce the missing arithmetic definition.
Any extension of the rational construction must additionally
control the new field-valued point coordinates and arithmetic
witnesses, or supply an independent existential definition
of the embedded rational field.  Under the claimed rational
definition, such rational isolation would allow an integral
basis construction of $\mathcal O_K$, but that additional
existential assertion has not been established here for
every number field.
